% 印刷xiv–xviii／物理15–19，原前言内的导论，不另造作者前言.
\subsection*{最优运输导论}\phantomsection\label{sec:introduction}\addcontentsline{toc}{subsection}{最优运输导论}
最优运输的历史始于很久以前的法国. 法国大革命前几年，Gaspard Monge在提交给科学院的一份报告\cite{b239}中提出如下问题\footnote{这发生在1781年；这里把他的问题改写为现代数学语言.}：给定$\mathbb R^d$上的两个质量密度$f,g\geq0$，满足$\int f(x)\,dx=\int g(y)\,dy=1$，寻找映射$\mathrm{T}:\mathbb R^d\to\mathbb R^d$，把第一个密度推到第二个密度上，即满足
\begin{equation}\tag{1}\label{eq:intro1}
\int_A g(y)\,dy=\int_{\mathrm{T}^{-1}(A)}f(x)\,dx\qquad\text{对任意Borel子集 }A\subset\mathbb R^d,
\end{equation}
并在所有满足此条件的映射中使
\[
M(\mathrm{T}):=\int_{\mathbb R^d}|\mathrm{T}(x)-x|f(x)\,dx
\]
最小. 这意味着我们有一群依照密度$f$分布于$\mathbb R^d$中的粒子，必须移动它们，使其形成预先指定密度为$g$的新分布. 移动方式应使平均位移最小. 在Monge的描述中，初始密度$f$表示要运到目标形状$g$的沙土分布. 这两个形状对应法语中的\textit{déblais}和\textit{remblais}，即挖方与填方. 显然，当时空间维数只考虑$d=2$或$3$. 映射$\mathrm{T}$描述我们必须最优选择的移动，$\mathrm{T}(x)$表示原来位于$x$处的粒子的目的地.

下文经常使用$X$上的测度$\mu$经可测映射$\mathrm{T}:X\to Y$得到的像测度；在问题最一般的表述中，测度将取代密度$f$与$g$. 这是$Y$上的测度，记为$\mathrm{T}_\#\mu$，与\eqref{eq:intro1}一样，由下式刻画：
\begin{align*}
(\mathrm{T}_\#\mu)(A)&=\mu(\mathrm{T}^{-1}(A))&&\text{对每个可测集 }A,\\
\text{或}\qquad\int_Y\phi\,d(\mathrm{T}_\#\mu)&=\int_X(\phi\circ \mathrm{T})\,d\mu&&\text{对每个可测函数 }\phi.
\end{align*}
更一般地，对于一般运输成本$c:X\times Y\to\mathbb R$，可以考虑
\[
(\mathrm{MP})\qquad\min\left\{M(\mathrm{T}):=\int c(x,\mathrm{T}(x))\,d\mu(x):\mathrm{T}_\#\mu=\nu\right\}.
\]
在欧氏框架内，如果两个测度$\mu,\nu$由密度$f,g$诱导，并假设$f,g,\mathrm{T}$足够正则且$\mathrm{T}$为单射，仅用换元公式便容易把等式$\nu=\mathrm{T}_\#\mu$变为偏微分方程
\begin{equation}\tag{2}\label{eq:intro2}\index{D039@Jacobi行列式（determinant of the Jacobian）|hyperpage}
g(\mathrm{T}(x))\det(D\mathrm{T}(x))=f(x).
\end{equation}
然而，这个方程关于$\mathrm{T}$高度非线性，这也是难以轻易分析Monge问题的困难之一. 例如，怎样证明极小元存在？通常的做法如下：取一列极小化序列$\mathrm{T}_n$，找到它的界，从而在某种拓扑中获得紧性；这里如果$\nu$的支集紧，则$\mathrm{T}_n$取值于共同的有界集$\operatorname{spt}(\nu)$，因而可以得到$\mathrm{T}_n$关于$L^\infty$弱星收敛的紧性. 随后取极限$\mathrm{T}_n\rightharpoonup \mathrm{T}$，并证明$\mathrm{T}$为极小元. 这要求泛函$M$对该收敛半连续；许多情形确实如此，例如$c$关于第二个变量凸. 我们需要
\[
\mathrm{T}_n\rightharpoonup \mathrm{T}\quad\Longrightarrow\quad\liminf_n M(\mathrm{T}_n)\geq M(\mathrm{T}),
\]
但还需要极限$\mathrm{T}$仍满足约束. 然而，偏微分方程的非线性使我们在只有弱收敛时无法证明这种稳定性. 读者可以尝试寻找一个弱收敛序列而其相应像测度不收敛的例子；这正是书末的\hyperref[ex:1]{习题(1)}.

Monge在\cite{b239}中分析了解的精细几何性质，强调了若干重要想法，我们将在\hyperref[sec:3]{第3章}看到：运输射线不相交、它们与某一族曲面正交，以及沿运输射线将各点单调排列是一种自然选择. 不过，他并没有真正解决问题，甚至没有处理极小元存在性. 此后150年，最优运输问题基本仍与法国紧密相连，科学院也为此设奖. 首个奖项由P. Appell凭一篇长篇论文\cite{b21}获得，该文改进了一些方面，但远未令人满意，也未处理存在性问题\footnote{关于Monge的工作，以及为该奖提交的论文内容等历史问题，读者可参见\cite{b181}；抱歉，该文为法语.}.

Monge问题一直没有解决：极小元是否存在，怎样刻画它？直到1940年代才取得进展. Kantorovich在1942年的工作，见\cite{b200}，才把它置于适当框架中，使人们能够着手研究，后来给出解并分析它们. 此后问题得到广泛推广，用很一般的成本函数$c(x,y)$取代欧氏距离$|x-y|$，并采用更一般的测度与空间. Kantorovich的主要思想是把Monge问题与线性规划联系起来. 他改变视角，用$X\times Y$上的测度$\gamma$描述粒子运动，要求$(\pi_x)_\#\gamma=\mu$且$(\pi_y)_\#\gamma=\nu$. 这些定义在$X\times Y$上的概率测度提供了另一种描述$\mu$中粒子位移的方式：不再给出每个原来位于$x$处的粒子的目的地$\mathrm{T}(x)$，而是对每一对$(x,y)$，给出从$x$前往$y$的粒子数量. 显然，这允许更一般的移动，因为原则上同一点$x$的粒子可以前往不同的目的地$y$. 若出现多个目的地，这种移动便无法用映射$\mathrm{T}$描述. 要极小化的成本变为简单的$\int_{X\times Y}c\,d\gamma$. 现在我们得到一个带有线性约束的线性问题，可以证明解存在，并运用对偶等凸优化技术刻画最优解，见\autoref{sec:1}.

某些情形下，尤其是$c(x,y)=|x-y|^2$时，甚至可以证明最优$\gamma$不允许这种质量分裂. 二次成本也非常自然，并因与动能的联系而在物理建模中广泛应用. 位于$x$的粒子只被送往唯一目的地$\mathrm{T}(x)$，由此给出原始Monge问题的解. Brenier在\cite{b82}中完成了这一点，还证明最优映射具有十分特殊的形式：$\mathrm{T}(x)=\nabla u(x)$，其中$u$是凸函数. 顺便说，这使问题与Monge--Ampère方程紧密联系：由\eqref{eq:intro2}得到
\[
\det(D^2u(x))=\frac{f(x)}{g(\nabla u(x))}.
\]
这是一个退化、非线性的方程，恰好在凸函数类中为椭圆型. Brenier还利用该结果给出向量映射的独特极分解定理，见\autoref{sec:1.7.2}：向量场可以写成凸函数梯度与保测映射的复合. 这推广了矩阵可以写成对称正定矩阵与正交矩阵之积这一事实.

Brenier的结果容易推广到其他关于差$x-y$严格凸的成本，也已推广到黎曼流形上的距离平方，见\cite{b231}. 但是，Monge最初提出的成本，即距离本身，要困难得多.

继法国学派之后，轮到俄罗斯数学家登场了. 从Kantorovich引入的精确方法出发，Sudakov\cite{b290}为原始Monge问题(MP)提出了一种解法. 成本$|x-y|$的Kantorovich问题的最优运输方案$\gamma$具有如下性质：$\mathbb R^d$可以分解为本质互不相交的线段之并，这些线段被$\gamma$保持，即$\gamma$集中于属于同一线段的点对$(x,y)$. 这些线段由一个Lipschitz函数构造，其水平集就是Monge“预见”的那些曲面. 随后，只须把问题化为一族一维问题. 如果$\mu\ll\mathscr L^d$，那么$\mu$在每条线段上诱导的测度也应绝对连续、没有原子. 而在一维，只要源测度没有原子，就可以定义单调递增运输，它对关于差$x-y$的任意凸成本都是最优的.

策略很清楚，但存在一个小缺陷：沿线段对$\mu$作分解所得测度的绝对连续性，听起来像Fubini型定理，却对任意线段族并不成立. 线段方向需要某种正则性. Ambrosio注意到这一点，并在\cite{b8,b10}中修复了它. 与此同时，Evans--Gangbo\cite{b159}通过一种与\hyperref[sec:4]{第4章}所见方法相关的方法获得了其他证明，可惜要求密度满足更强假设；Caffarelli--Feldman--McCann\cite{b101}和独立工作的Trudinger--Wang\cite{b291}则通过严格凸成本逼近得到了证明.

\phantomsection\label{idx:i123-pxvii}\index{Q122@二次成本（quadratic cost）!123@流形上（on manifolds）|hyperpage}
在为最优映射存在性付出许多努力后，人们也研究了其正则性. 这一方面的主要参考是Caffarelli的工作，他通过研究上述Monge--Ampère方程，证明了二次情形的正则性. 令人惊讶的是，本世纪初，Ma--Trudinger--Wang\cite{b219}找到了不同成本下正则性的关键. 特别是，他们对$\mathbb R^d$上的$C^4$成本找到了保证正则性的条件，即关于四阶导数的一些不等式. 该条件可以推广到光滑流形上的距离平方，此时假设变为底流形的曲率条件. Loeper后来在\cite{b216}中证明这些条件是尖锐的. 正则性是一个优美而精细的主题，本书无法给予它本应获得的全部篇幅；更多细节和参考见\autoref{sec:1.7.6}.

最优运输理论不能仅归结为最优映射的存在性和性质. 这一理论的成功，与它和许多其他数学分支的联系密不可分. 部分联系通过最优映射实现，例如一些可以借助它证明或重新证明的几何与泛函不等式. 本书只介绍等周不等式，见\hyperref[sec:2.5.3]{第2.5.3节}，以及Brunn--Minkowski不等式，见\hyperref[sec:7.4.2]{第7.4.2节}. 应当强调，定量等周不等式中最精细的进展之一来自Figalli--Maggi--Pratelli，其证明强烈依赖最优运输工具\cite{b167}.

另一方面，许多最优运输应用通过其定义的距离实现，即Wasserstein距离，见\hyperref[sec:5]{第5章}. 对每一对测度，可以根据成本$|x-y|^p$下将$\mu$运往$\nu$的最小成本，关联一个记为$W_p(\mu,\nu)$的量. 可以证明它是距离，并将概率测度的弱收敛度量化，至少在紧空间上如此. 这一距离在研究某些偏微分方程时很有用. 一些演化方程，尤其是可能带非局部项的抛物型方程，可以解释为这一距离的梯度流，即最速下降曲线，见\hyperref[sec:8]{第8章}. 这一思想最初由\cite{b198,b246}强调，可用来证明解存在或逼近解. 对另一些偏微分方程，沿两个解对Wasserstein距离求导，可以成为证明稳定性、唯一性或向稳态收敛速率的技术工具，见\hyperref[sec:5.3.5]{第5.3.5节}. 最后，还有一些演化模型与涉及动能的作用量极小化有关——这在物理学中十分常见，因为密度曲线相对于$W_2$距离的速度，恰好是某种动能——或者与方程中出现的特殊凸函数的梯度有关，见\hyperref[sec:8.4.4]{第8.4.4节}.

配备距离$W_p$的概率测度空间$\mathbb W_p$的结构，过去和现在都受到许多关注. 尤其是其测地线，以及某些泛函沿这些测地线的凸性，既在参考书\cite{b15}发展的度量梯度流理论中发挥关键作用，也具有几何方面的推论. 在黎曼流形上，熵$E(\varrho):=\int\varrho\log\varrho$是否沿这些测地线凸，与Ricci曲率的下界等价. 由此产生了度量测度空间上广泛的分析理论，其中将这种凸性作为曲率下界的定义，见\cite{b218,b288,b289}. 近年来，这一理论因Ambrosio、Gigli、Savaré、Kuwada、Ohta及其合作者的许多新结果而大有进展，例如\cite{b18,b183}.

从建模角度看，最优运输可出现在与空间经济学、交通、网络和集体运动等有关的许多领域，但至少对某些具体应用，Monge--Kantorovich模型是否合适，仍可受到质疑. 这促使人们研究另一类模型：交通拥挤对应更“凸”的模型，组织高效运输网络对应更“凹”的模型. 这些模型通常写成带散度约束的流极小化形式，在某种意义上是原始Monge成本的变体. 实际上，原始Monge问题，即成本为$|x-y|$时从$\mu$到$\nu$的最优运输，等价于在$\nabla\cdot w=\mu-\nu$的条件下极小化向量场$w$的$L^1$范数. \hyperref[sec:4]{第4章}介绍这一问题，\hyperref[sec:4.4]{第4.4节}则把交通拥挤与分支运输模型作为其变体介绍.

最后，特别是由于图像处理中的应用，见\hyperref[sec:2.5.1]{第2.5.1节}和\hyperref[sec:5.5.5]{第5.5.5节}，高效计算或逼近两个测度之间的最优运输及Wasserstein距离，近来变得至关重要. 这是一个新兴而活跃的研究领域；\hyperref[sec:6]{第6章}介绍的仅是若干经典方法. 本书无意在这方面穷尽全部内容，只希望对这些主题提供一些说明.

我们无意在最优运输的“纯粹”和“应用”两个侧面之间筑墙. 近年来，两者都以令人瞩目的速度进步. 本书致力于那些可能更能吸引以下读者的理论主题：关注经济学、图像处理、社会和生物演化现象、流体力学中的建模问题，关注偏微分方程应用，或关注数值问题的读者. 无论这篇简短导论，还是整本书，都不可能完全概括这些主题正在探索的新方向.

我们只希望让读者体会这些主题的趣味，并严格分析其背后的数学.
